\scanpage{76}
% archive foliation: 75
existe et est réunion des $\varphi_{S'/S}(\mathfrak X_i)$.

Démonstration. Il suffit de prouver ii), et pour cela de prouver
que l'hyp. faite implique la condition de recollement. Or,
si $v:T'\to\mathfrak X$ est donné, on introduit les
$T'_i=v^{-1}(\mathfrak X_i)$, \deleted{$\unsure{-}$} et les
\[
T_i=\complement\varphi_T(T'-T'_i),
\]
qui sont des ouverts de $T$, on doit prouver que $T=\bigcup T_i$;
\deleted{et pour ceci, il}
Or $t\in T_i$ signifie que la restriction à $T'_t$ de $\unsure{v|\varphi}$ applique
$T'_t$ dans $\mathfrak X_i$ : donc une section de
$(\mathfrak X_s/S'_s)_{/k(s)}$ définie sur l'extension $k(t)$ de $k(s)$,
doit être contenue dans un des $\mathfrak X_i$.
\margin{\diagram{\begin{tikzcd}[column sep=small,row sep=small]
& \mathfrak X \arrow[d] \\
S & S' \arrow[l,"\varphi"'] \\
T \arrow[u] & T' \arrow[l,"\varphi_T"'] \arrow[u,no head] \arrow[uu,bend right=35,dashed,"\unsure{v|u}"']
\end{tikzcd}}}
\margin{\diagram{$\begin{array}{cc}
&\mathfrak X_s\\
k(s)&S'_s\\[1em]
k(t)&T'_t
\end{array}$}}
\deleted{Or \unsure{si on décompose} $S'_s$
\inserted{en des $S_s'^\alpha$}
\unsure{suivant les comp. connexes}
\inserted{$S_s'^\alpha$} $S_s'^\alpha$,
\unsure{on est ramené à} \unsure{ce que chaque} $\alpha$
\unsure{une section de} $\mathfrak X_i^\alpha$ \unsure{sur}}

\uline{Corollaire 1} Si Toute partie \deleted{\unsure{tt}} finie
\inserted{ayant $\leq n(s)$ éléments} d'un $\mathfrak X_s$ est contenue
dans un ouvert affine de $\mathfrak X$, \deleted{alors}
[en particulier si \deleted{$\mathfrak X_s$} \inserted{$\mathfrak X/S$} est
quasi-projectif] alors $\varphi_{S'/S}(\mathfrak X)$ existe, et est
\margin{Soit $n(s)$ le \deleted{\illegible{2}} \unsure{nombre}
\unsure{géom.} \unsure{de points} de $S'/S$}
